\chapter{Uogólniona wymierna kontrakcja i łańcuchy \texorpdfstring{$C_{p,q}$}{Cpq}}
\label{ch:tg-rb-general}

W poprzednim rozdziale rozpatrzyliśmy konfigurację $C_p$ opisaną przez
łańcuch sfer o wagach $-(p+2),-2,\ldots,-2$. W niniejszym
rozdziale wyprowadzimy uogólnienie $C_{p,q}$, gdzie liczby $p>q\geq1$
są względnie pierwsze. Podstawowym narzędziem będzie \emph{ujemne
rozwinięcie ułamka łańcuchowego}. Odróżniamy przy tym rachunki
algebraiczne, które wykonamy w całości, od geometrycznego twierdzenia
o istnieniu wymiernej kuli homologii.

\section{Ujemne ułamki łańcuchowe: definicja i jednoznaczność}
\label{sec:tg-general-cf}

\begin{definicja}
Dla liczb całkowitych $a_1,\ldots,a_r$, gdzie $r\geq1$ oraz
$a_i\geq2$, definiujemy \emph{ujemny ułamek łańcuchowy}
rekurencyjnie przez $[a_r]^-:=a_r$ oraz
$[a_i,\ldots,a_r]^-:=a_i-1/[a_{i+1},\ldots,a_r]^-$.
\end{definicja}

Nie jest to zwykłe dodatnie rozwinięcie Euklidesa. W szczególności
w każdym kroku odejmujemy \emph{odwrotność} kolejnego ogona.

\begin{twierdzenie}[Algorytm ujemnego rozwinięcia]
\label{thm:tg-cf-algorithm}
Niech $A,B\in\mathbb N$ spełniają $A>B\geq1$ i
$\gcd(A,B)=1$. Istnieje dokładnie jeden skończony ciąg
$a_1,\ldots,a_r\geq2$ taki, że $A/B=[a_1,\ldots,a_r]^-$. 
Otrzymujemy go, biorąc $a_1=\lceil A/B\rceil$, a jeśli
$B>1$, zastępując $A/B$ ułamkiem $B/(a_1B-A)$ i powtarzając
postępowanie. Proces kończy się, gdy mianownik wynosi jeden.
\end{twierdzenie}

\begin{proof}
Jeżeli $B=1$, jedynym rozwinięciem jest $[A]^-$: każde
rozwinięcie o co najmniej dwóch wyrazach ma wartość niecałkowitą,
ponieważ jego pierwszy ogon jest większy od jedności. Niech
$B>1$. Ponieważ $A/B$ jest nieskracalny, nie jest całkowity.
Definiujemy $a_1:=\lceil A/B\rceil\geq2$ oraz
$B':=a_1B-A$. Wówczas $0<B'<B$ i $\gcd(B,B')=1$.
Zatem
$A/B=a_1-B'/B=a_1-1/(B/B')$.
Dalszy ułamek ma mianownik $B'$ mniejszy od $B$, toteż
indukcja kończy konstrukcję.

Dla dowodu jednoznaczności zauważmy, że każdy ogon
$[a_i,\ldots,a_r]^->$ jest większy od jedności. Stąd
$a_1-1<[a_1,\ldots,a_r]^-\leq a_1$, przy czym równość
po prawej zachodzi wyłącznie dla $r=1$.
Pierwszy wyraz jest więc wymuszony jako sufit danego ułamka;
po jego usunięciu wymuszony jest także następny ogon.
Indukcja daje jednoznaczność wszystkich wyrazów.
\end{proof}

\begin{przyklad}[Trzy kolejne obliczenia]
Dla $(p,q)=(5,1)$ zachodzi $25/4=[7,2,2,2]^-$.
Dla $(p,q)=(5,2)$ obliczamy
$25/9=3-2/9=3-1/(9/2)=3-1/(5-1/2)$, czyli
$25/9=[3,5,2]^-$. Dla $(p,q)=(5,3)$ mamy
$25/14=2-3/14=2-1/(14/3)=2-1/(5-1/3)$,
zatem $25/14=[2,5,3]^-$. Kolejność wag ma znaczenie dla
konwencji orientacji i oznaczania generatorów brzegu.
\end{przyklad}

\section{Definicja \texorpdfstring{$C_{p,q}$}{Cpq} i forma przecięcia}
\label{sec:tg-general-config}

\begin{definicja}[Konfiguracja $C_{p,q}$]
Niech $p,q\in\mathbb N$, $p>q\geq1$, $\gcd(p,q)=1$, i niech
$p^2/(pq-1)=[a_1,\ldots,a_r]^-$, $a_i\geq2$.
Przez $C_{p,q}$ oznaczamy zwartą, spójną, zorientowaną,
gładką czterorozmaitość otrzymaną przez liniowe plumbingowanie
wiązek dyskowych nad $S^2$ o liczbach Eulera
$-a_1,\ldots,-a_r$, przy czym każda para sąsiednich sfer
przecina się poprzecznie i dodatnio w jednym punkcie.
Forma przecięcia $Q_{p,q}$ w bazie klas tych sfer ma macierz
\[
Q_{p,q}=\begin{pmatrix}
-a_1&1&0&\cdots&0\\
1&-a_2&1&\ddots&\vdots\\
0&1&\ddots&\ddots&0\\
\vdots&\ddots&\ddots&-a_{r-1}&1\\
0&\cdots&0&1&-a_r
\end{pmatrix}
\]
(dla $r=1$ jest to macierz $[-a_1]$).
\end{definicja}

\begin{center}
\begin{tikzpicture}[>=Stealth,scale=.9]
  \node[font=\small\bfseries,text=tgNavy] at (-3.7,.92) {$C_{5,1}$};
  \node[tgRedVertex] (a) at (-5.2,0) {$-7$};
  \node[tgBlueVertex] (b) at (-4.2,0) {$-2$};
  \node[tgBlueVertex] (c) at (-3.2,0) {$-2$};
  \node[tgBlueVertex] (d) at (-2.2,0) {$-2$};
  \draw[tgTree] (a)--(b)--(c)--(d);
  \node[font=\small\bfseries,text=tgNavy] at (2.4,.92) {$C_{5,2}$};
  \node[tgRedVertex] (u) at (1,0) {$-3$};
  \node[tgBlueVertex] (v) at (2.3,0) {$-5$};
  \node[tgBlueVertex] (w) at (3.6,0) {$-2$};
  \draw[tgTree] (u)--(v)--(w);
  \node[font=\scriptsize,text=tgNavy] at (-3.7,-.68) {$25/4=[7,2,2,2]^-$};
  \node[font=\scriptsize,text=tgNavy] at (2.35,-.68) {$25/9=[3,5,2]^-$};
\end{tikzpicture}
\end{center}

\begin{lemat}[Wyznaczniki kolejnych ogonów]
\label{lem:tg-general-minors}
Niech $a_1,\ldots,a_r\geq2$. Oznaczmy przez $D_i$ wyznacznik
macierzy trójdiagonalnej o przekątnej
$a_i,\ldots,a_r$ i elementach $-1$ bezpośrednio przy przekątnej,
a ponadto przyjmijmy $D_{r+1}=1$, $D_{r+2}=0$. Wtedy
$D_i=a_iD_{i+1}-D_{i+2}$ i
$[a_i,\ldots,a_r]^-=D_i/D_{i+1}$. Wszystkie $D_i$ dla
$1\leq i\leq r+1$ są dodatnie, zaś $D_i>D_{i+1}$ dla
$1\leq i\leq r$.
\end{lemat}

\begin{proof}
Rozwinięcie wyznacznika względem pierwszego wiersza daje
$D_i=a_iD_{i+1}-D_{i+2}$. Dla $i=r$ mamy $D_r=a_r$ i
$D_{r+1}=1$. Indukcja wsteczna prowadzi do
$D_i\geq2D_{i+1}-D_{i+2}>D_{i+1}>0$.
Dzieląc rekurencję przez $D_{i+1}$ otrzymujemy
$D_i/D_{i+1}=a_i-D_{i+2}/D_{i+1}$, co jest dokładnie
rekurencją definiującą ujemny ułamek łańcuchowy.
\end{proof}

\begin{twierdzenie}[Ujemna określoność, wyznacznik i homologia]
\label{thm:tg-general-algebra}
Dla dowolnych względnie pierwszych $p>q\geq1$ macierz
$Q_{p,q}$ jest ujemnie określona i
$\det Q_{p,q}=(-1)^r p^2$.
Ponadto jej cokernel jest cykliczny i
\[
H_1(\partial C_{p,q};\mathbb Z)
\cong\operatorname{coker}Q_{p,q}
\cong\mathbb Z/p^2\mathbb Z.
\]
W szczególności $\partial C_{p,q}$ jest wymierną sferą
homologiczną.
\end{twierdzenie}

\begin{proof}
Niech $A:=-Q_{p,q}$. Z poprzedniego lematu mamy
$\det A=D_1=p^2$, gdyż
$D_1/D_2=p^2/(pq-1)$ jest ułamkiem nieskracalnym,
a $D_1,D_2$ są całkowite i względnie pierwsze
(ten ostatni fakt wynika rekurencyjnie z
$\gcd(D_i,D_{i+1})=\gcd(D_{i+1},D_{i+2})=1$).
W szczególności $D_2=pq-1$.

Dowód określoności nie wymaga wyliczania wartości własnych.
Dla $r\geq2$ i $x=(x_1,\ldots,x_r)\in\mathbb R^r$
zachodzi tożsamość
\[
x^TAx=\sum_{i=1}^{r-1}(x_i-x_{i+1})^2
+(a_1-1)x_1^2+(a_r-1)x_r^2
+\sum_{i=2}^{r-1}(a_i-2)x_i^2.
\]
Każdy składnik po prawej jest nieujemny, a pierwszy i ostatni
składnik brzegowy wymuszają $x_1=x_r=0$ przy zerowej sumie;
różnice sąsiednich współrzędnych wymuszają wtedy $x=0$.
Zatem $A$ jest dodatnio określona. Dla $r=1$ mamy $A=[a_1]$.
Otrzymujemy ujemną określoność $Q_{p,q}$ i wzór na wyznacznik.

Cokernel $Q_{p,q}$ jest generowany przez klasy
$e_1,\ldots,e_r$ z relacjami
$-a_1e_1+e_2=0$, $e_{i-1}-a_ie_i+e_{i+1}=0$
oraz $e_{r-1}-a_re_r=0$ (dla $r=1$ tylko
$-a_1e_1=0$). Pierwsza relacja wyraża $e_2$ przez $e_1$,
a każda następna wyraża $e_{i+1}$ przez poprzednie.
W konsekwencji grupa jest generowana przez $e_1$.
Ponieważ macierz jest nieosobliwa, grupa jest skończona rzędu
$|\det Q_{p,q}|=p^2$, więc musi być cykliczna rzędu $p^2$.
Identyfikację z $H_1$ brzegu uzyskaliśmy wcześniej z
ciągu dokładnego pary i dualności Poincarégo--Lefschetza.
\end{proof}

\begin{przyklad}[Rachunek dla $C_{5,2}$]
Dla $25/9=[3,5,2]^-$ otrzymujemy
\[
Q_{5,2}=\begin{pmatrix}-3&1&0\\1&-5&1\\0&1&-2\end{pmatrix},
\qquad\det Q_{5,2}=-25.
\]
Relacje $-3e_1+e_2=0$ oraz
$e_1-5e_2+e_3=0$ dają $e_2=3e_1$ i $e_3=14e_1$.
Ostatnia relacja $e_2-2e_3=0$ staje się
$-25e_1=0$. W ten sposób bez użycia postaci Smitha
obliczamy $H_1(\partial C_{5,2};\mathbb Z)\cong\mathbb Z/25\mathbb Z$.
\end{przyklad}

\section{Identyfikacja brzegu i forma splecenia}
\label{sec:tg-general-boundary}

\begin{twierdzenie}[Brzeg konfiguracji $C_{p,q}$]
\label{thm:tg-general-lens}
Przy konwencji $L(n,s)=S^3_{n/s}(U)$ dla
niezorientowanego węzła trywialnego $U\subset S^3$,
z orientacją wyznaczoną przez chirurgię, istnieje
orientacyjny dyfeomorfizm
$\partial C_{p,q}\cong L(p^2,1-pq)$.
\end{twierdzenie}

\begin{proof}
Diagram Kirby'ego $C_{p,q}$ jest liniowym łańcuchem
ramowanych okręgów trywialnych o współczynnikach
$-a_1,\ldots,-a_r$ i liczbach splecenia sąsiednich
składowych równych $+1$.
Stosując kolejno geometryczną regułę \emph{slam-dunk}
(udowodnioną wcześniej dla usuwania sąsiedniej składowej),
redukujemy diagram do pojedynczej składowej z wymiernym
współczynnikiem
$-[a_1,\ldots,a_r]^-= -p^2/(pq-1)=p^2/(1-pq)$.
Jest to chirurgiczne przedstawienie $L(p^2,1-pq)$.
Argument nie korzysta wyłącznie z wyznacznika macierzy:
niezależnie ustala klasę krzywej południkowej podczas sklejenia.
\end{proof}

\begin{stwierdzenie}[Generator, forma splecenia i metabolizer]
\label{prop:tg-general-linking}
Niech $e_1$ oznacza klasę pierwszego wierzchołka w
$H_1(\partial C_{p,q};\mathbb Z)=\operatorname{coker}Q_{p,q}$.
Przy konwencji, w której formę splecenia brzegu reprezentuje
$-Q_{p,q}^{-1}$, zachodzi
\[
\lambda(e_1,e_1)=\frac{pq-1}{p^2}\pmod{\mathbb Z}.
\]
Jedyna podgrupa rzędu $p$, czyli $p\langle e_1\rangle$,
jest równa swojemu dopełnieniu ortogonalnemu względem
$\lambda$. Jest więc \emph{metabolizerem} tej formy.
\end{stwierdzenie}

\begin{proof}
Wyznacznik macierzy powstałej po usunięciu pierwszego
wiersza i kolumny $Q_{p,q}$ jest równy
$(-1)^{r-1}D_2$, więc z reguły dopełnień algebraicznych
\[
-(Q_{p,q}^{-1})_{11}=
-\frac{(-1)^{r-1}D_2}{(-1)^rD_1}
=\frac{D_2}{D_1}=
\frac{pq-1}{p^2}.
\]
Dla $r=1$ przyjmujemy wyznacznik macierzy pustej równy jeden.
Ponieważ $e_1$ ma rząd $p^2$ i $\gcd(p,pq-1)=1$,
forma jest niezdegenerowana. Dla generatora $pe_1$ podgrupy
rzędu $p$ mamy
$\lambda(pe_1,pe_1)=pq-1\equiv0\pmod{\mathbb Z}$.
Cała podgrupa jest zatem izotropowa.
Jej dopełnienie ortogonalne ma rząd $p^2/p=p$,
więc z zawierania i zgodności rzędów otrzymujemy równość.
\end{proof}

\begin{uwaga}[Konwencje znaku]
W literaturze często definiuje się $L(n,s)$ przez
\emph{dodatni} ułamek chirurgii $n/s$, a niektóre teksty
przyjmują przeciwną orientację brzegu plumbingowania.
Wówczas można spotkać $L(p^2,pq-1)$ zamiast
$L(p^2,1-pq)$. Przy korzystaniu z takich wyników
należy jawnie odwrócić orientację odpowiedniego brzegu,
a nie tylko zmienić znak w liczniku lub mianowniku.
\end{uwaga}

\begin{przyklad}[Dokładna kontrola konfiguracji $C_{7,3}$]
\label{ex:tg-general-C73}
Dla $(p,q)=(7,3)$ mamy $p^2=49$ oraz
$pq-1=20$. Algorytm ujemnego rozwinięcia daje
\[
\frac{49}{20}=[3,2,6,2]^-.
\]
Sprawdzamy to od końca:
$[6,2]^-=11/2$,
$[2,6,2]^-=2-2/11=20/11$,
a zatem $[3,2,6,2]^-=3-11/20=49/20$.
Macierz przecięcia konfiguracji jest równa
\[
Q_{7,3}=
\begin{pmatrix}
-3&1&0&0\\
1&-2&1&0\\
0&1&-6&1\\
0&0&1&-2
\end{pmatrix}.
\]
Od końca wyznaczniki macierzy ogonowych
dla macierzy $-Q_{7,3}$
wynoszą odpowiednio
$D_5=1$, $D_4=2$, $D_3=11$,
$D_2=20$ oraz $D_1=49$.
Stąd $\det Q_{7,3}=49$ i
\[
H_1(\partial C_{7,3};\mathbb Z)
\cong\mathbb Z/49\mathbb Z,\qquad
\lambda(e_1,e_1)=\frac{20}{49}\pmod{\mathbb Z}.
\]
Dla elementu $7e_1$ otrzymujemy
$\lambda(7e_1,7e_1)=20\equiv0\pmod{\mathbb Z}$.
Podgrupa $\langle7e_1\rangle$ rzędu $7$
jest zatem metabolizerem.
Przy ustalonej konwencji brzeg jest
$L(49,-20)$. Wymierna kontrakcja
takiej konfiguracji, o ile jest ona
gładko zanurzona w wyjściowej
czterorozmaitości, zwiększa jej sygnaturę
o $4$ i zmniejsza charakterystykę Eulera o $4$.
Żaden z tych rachunków nie dowodzi
istnienia konkretnego zanurzenia konfiguracji.
\end{przyklad}

\section{Wymierna kula homologii i zmiana niezmienników}
\label{sec:tg-general-operation}

\begin{twierdzenie}[Istnienie $B_{p,q}$; wynik przyjmowany]
\label{thm:tg-general-ball}
Dla każdego $p>q\geq1$ ze $\gcd(p,q)=1$ istnieje
zwarta, spójna, zorientowana, gładka czterorozmaitość
$B_{p,q}$ o brzegu orientacyjnie dyfeomorficznym z
$\partial C_{p,q}$, dla której
\[
H_*(B_{p,q};\mathbb Q)\cong H_*(B^4;\mathbb Q),
\qquad H_1(B_{p,q};\mathbb Z)\cong\mathbb Z/p\mathbb Z.
\]
Można wybrać standardowy model, w którym
$\pi_1(B_{p,q})\cong\mathbb Z/p\mathbb Z$, a mapa
$\pi_1(\partial B_{p,q})\to\pi_1(B_{p,q})$ jest surjektywna.
\end{twierdzenie}

\begin{uwaga}
Istnienie $B_{p,q}$ z tym brzegiem jest nietrywialnym
twierdzeniem geometrycznym, wynikającym z konstrukcji
wymiernych kul homologii (Casson--Harer i ich uogólnienia).
Powyższych faktów nie da się wyprowadzić z samej równości
$|H_1(\partial C_{p,q})|=p^2$ ani z istnienia metabolizera.
Traktujemy je tu jako wynik zewnętrzny, wskazując jego dokładną
rolę w dalszym rozumowaniu.
\end{uwaga}

\begin{definicja}[Uogólniona wymierna kontrakcja]
Niech $X$ będzie zamkniętą, spójną, zorientowaną,
gładką czterorozmaitością zawierającą gładko zanurzoną
kopię $C_{p,q}$ jako podrozmaitość z brzegiem.
Kładziemy $W:=X\setminus\operatorname{int}C_{p,q}$.
Dla ustalonego dyfeomorfizmu
$\varphi:\partial B_{p,q}\to\partial C_{p,q}$ zachowującego
orientację tych dwóch brzegów definiujemy
$X_{p,q}:=W\cup_{\varphi}B_{p,q}$.
Ponieważ $\partial W=-\partial C_{p,q}$, jest to
sklejenie odwracające orientację między $\partial W$
a $\partial B_{p,q}$. Wynik nazywamy
\emph{uogólnioną wymierną kontrakcją} $X$ wzdłuż $C_{p,q}$.
\end{definicja}

\begin{twierdzenie}[Obliczenie niezmienników]
\label{thm:tg-general-numeric}
Niech $r$ będzie długością ujemnego rozwinięcia
$p^2/(pq-1)$, a $X_{p,q}$ powstanie z zamkniętej,
zorientowanej, gładkiej czterorozmaitości $X$
przez uogólnioną wymierną kontrakcję wzdłuż $C_{p,q}$.
Wówczas
\[
\begin{aligned}
e(X_{p,q})&=e(X)-r,&
\sigma(X_{p,q})&=\sigma(X)+r,\\
b_2^+(X_{p,q})&=b_2^+(X),&
b_2^-(X_{p,q})&=b_2^-(X)-r.
\end{aligned}
\]
Stąd $\chi_h:=(e+\sigma)/4$ pozostaje bez zmian,
natomiast $c_{\mathrm{num}}:=2e+3\sigma$ wzrasta o $r$.
Przez $\chi_h$ rozumiemy tu jedynie liczbowe wyrażenie
$(e+\sigma)/4$; nie zakładamy istnienia struktury
zespolonej, która uzasadniałaby interpretację holomorficzną.
\end{twierdzenie}

\begin{proof}
Brzeg $Y=\partial C_{p,q}$ jest wymierną sferą homologiczną.
Ciągi Mayera--Vietorisa nad $\mathbb Q$ dają rozkłady
$H_2(X;\mathbb Q)\cong H_2(W;\mathbb Q)\oplus
H_2(C_{p,q};\mathbb Q)$ oraz
$H_2(X_{p,q};\mathbb Q)\cong H_2(W;\mathbb Q)$.
Sumę pierwszą można przedstawić przez powierzchnie o nośnikach
w rozłącznych częściach, zatem jest ona ortogonalna względem
formy przecięcia. Ponieważ $Q_{p,q}$ jest ujemnie określona
rangi $r$, usunięcie konfiguracji zmniejsza $b_2^-$ o $r$
i nie zmienia $b_2^+$. Wynika stąd zmiana sygnatury o $r$.

Z rozkładu uchwytowego plumbingowania, który ma jedną
komórkę wymiaru zero i $r$ dwukomórek, otrzymujemy
$e(C_{p,q})=1+r$. Natomiast $e(B_{p,q})=1$
i $e(Y)=0$. Addytywność charakterystyki Eulera prowadzi do
$e(X_{p,q})-e(X)=1-(1+r)=-r$.
Pozostałe wzory są prostym podstawieniem.
\end{proof}

\begin{wniosek}[Prosta spójność po sklejeniu]
Jeżeli w sytuacji poprzedniego twierdzenia dopełnienie
$W=X\setminus\operatorname{int}C_{p,q}$ jest po prostu spójne,
a model $B_{p,q}$ wybieramy jak w twierdzeniu
\ref{thm:tg-general-ball}, to $X_{p,q}$ również jest po prostu
spójna.
\end{wniosek}

\begin{proof}
Z twierdzenia van Kampena otrzymujemy
$\pi_1(X_{p,q})\cong\pi_1(W)*_{\pi_1(Y)}\pi_1(B_{p,q})$.
Gdy $\pi_1(W)=1$, pierwsza grupa znika, a iloraz
$\pi_1(B_{p,q})$ przez normalne domknięcie obrazu
$\pi_1(Y)$ jest trywialny, ponieważ obraz jest całą grupą.
\end{proof}

\begin{przyklad}[Co wnosi parametr $q$?]
Dla $(p,q)=(5,1)$ mamy łańcuch długości $r=4$ o wagach
$-7,-2,-2,-2$. Dla $(p,q)=(5,2)$ łańcuch ma długość
$r=3$ i wagi $-3,-5,-2$. Oba brzegi mają tę samą grupę
$H_1\cong\mathbb Z/25\mathbb Z$, lecz odpowiednie formy
splecenia w kanonicznych generatorach mają wartości
odpowiednio $4/25$ i $9/25$; również współczynniki chirurgii
są różne. Jeśli konfiguracje występują jako podrozmaitości
w różnych czterorozmaitościach, operacja zmienia ich sygnatury
o odpowiednio $4$ i $3$. Z samych danych homologicznych
nie wolno wnioskować, że otrzymane rozmaitości są dyfeomorficzne.
\end{przyklad}

\begin{center}
\begin{tikzpicture}[>=Stealth]
  \node[draw=tgNavy,fill=tgBlue!9,rounded corners=2pt,
    minimum width=27mm,minimum height=11mm,align=center,font=\small] (w) at (0,1.6) {$W$};
  \node[draw=tgCoral,fill=tgCoral!13,rounded corners=3pt,
    minimum width=29mm,minimum height=13mm,align=center,font=\small] (c) at (-3,0) {$C_{p,q}$};
  \node[draw=tgTeal,fill=tgTeal!15,rounded corners=3pt,
    minimum width=29mm,minimum height=13mm,align=center,font=\small] (b) at (3,0) {$B_{p,q}$};
  \draw[tgAccent,-{Stealth[length=2.9mm]}] (w.south west)--(c.north east);
  \draw[tgTree,-{Stealth[length=2.9mm]}] (w.south east)--(b.north west);
  \node[font=\footnotesize,text=tgNavy] at (-3,-1.1) {$X=W\cup C_{p,q}$};
  \node[font=\footnotesize,text=tgNavy] at (3,-1.1) {$X_{p,q}=W\cup B_{p,q}$};
  \node[font=\scriptsize,text=tgNavy] at (0,-1.8)
    {$\partial C_{p,q}\cong\partial B_{p,q}\cong L(p^2,1-pq)$};
\end{tikzpicture}
\end{center}

\section{Granica metody i jej zastosowania}
\label{sec:tg-general-limit}

Wymierna kontrakcja zmienia formę przecięcia i liczby Betti'ego,
ale te dane nie wystarczają do dowodu egzotyczności.
W konstrukcjach gładkich trzeba jeszcze wykazać odpowiednią
prostą spójność i skorzystać z niezmienników zależnych od
struktury różniczkowej. W szczególności wykrywanie
nierównoważnych struktur gładkich wymaga teorii cechowania.
Jednym z jej podstawowych narzędzi są niezmienniki
Seiberga--Wittena, które wprowadzimy w kolejnym rozdziale.

\begin{uwaga}[Uogólniona kontrakcja a typ gładki]
Ciąg Mayera--Vietorisa i algebra macierzy przecięcia
kontrolują homologię, sygnaturę i część informacji
o formie splecenia. Nie dowodzą jednak istnienia
zanurzonej konfiguracji $C_{p,q}$ w zadanej
czterorozmaitości, ani dyfeomorficznego typu jej
dopełnienia. W zastosowaniach do egzotycznych struktur
gładkich każdy z tych punktów wymaga osobnego
argumentu geometrycznego.
\end{uwaga}

